\documentclass[11pt]{article}
\usepackage[T1]{fontenc}
\usepackage[utf8]{inputenc}
\usepackage{lmodern}
\usepackage{amsmath,amssymb,mathrsfs,booktabs,hyperref}
\usepackage[margin=23mm]{geometry}
\hypersetup{hidelinks,pdftitle={Yang--Mills: complete research proofs}}
\setlength{\parindent}{0pt}
\setlength{\parskip}{6pt}
\begin{document}
\section*{Global classical evolution of the original compactly supported family}

30 September 2026. This continuation solves the classical Cauchy problem
for every input of the retained moment-map family, through its actual
cutoff region. It applies the established global existence theorem of
Sung-Jin Oh, with the exact hypotheses and convention map below. The
bundle descent, parameter derivatives, support law, core comparison,
energy kernel and dilation formulas are proved here. This is a classical
solution family; construction of the physical quantum state remains the
programme goal.

\subsection*{1. Original data, without removing the cutoff}

Use the principal bundle, ordered line coordinates and representation in
\href{https://github.com/KokunoYumeto/yang-mills-interacting-workbench/blob/a15d9ffe8edaec69cdf91425dabbfe3170bd5ae9/yang-mills/continuations/20260930-s6-ns-moment-map-bridge/PROOF.md#7-the-global-triality-to-profile-morphism}{the moment-map proof, Theorem 7.1}.
The base comparison, dimensions and full parameter connection are in
\href{https://github.com/KokunoYumeto/yang-mills-interacting-workbench/blob/a15d9ffe8edaec69cdf91425dabbfe3170bd5ae9/yang-mills/continuations/20260930-s6-ns-moment-map-bridge/HIGHER_CARRIER_AND_EVOLUTION.md#2-the-reduction-carrier-and-its-exact-pullback}{the higher-carrier proof, Sections 2–4}.
In particular
\[
 X\cong S^6,\quad K=\rho(\operatorname{Sp}(1))\subset G=\operatorname{Spin}(8)
 \subset L=F_4,\quad D=L/K,\quad E_D=L\times_K W.
\]
The map \(r_H:X\to D\) pulls \(L\to D\) back to \(P_H\), and pulls \(E_D\)
back to the original \(\mathcal W_H\). Here \(\dim D=49\), while
\(\operatorname{rank}E_D=24\). The representation retains
\[
 v=(\zeta;\alpha,\beta;\gamma,\delta),\qquad
 c_1=\alpha\mathbf i\bar\alpha,\quad
 c_2=\beta\mathbf j\bar\beta,\quad
 c_3=\gamma\mathbf k\bar\gamma.
 \tag{CE1}
\]
The unused coordinates in this polynomial map are still part of \(W\).
Retain the fixed \(0<r<R\), the original real
\(\chi\in C_c^\infty(B_R(0))\), and \(\chi=1\) on \(B_r(0)\).
With \(\chi_i=\partial_i\chi\), the three original curls are
\[
 \begin{aligned}
 w^{(1)}&=(\chi+x_2\chi_2,\,-x_2\chi_1,\,0),\\
 w^{(2)}&=(-x_1\chi_2,\,\chi+x_1\chi_1,\,0),\\
 w^{(3)}&=(-x_1\chi_3,\,0,\,\chi+x_1\chi_1).
 \end{aligned}\tag{CE2}
\]
Thus
\[
 C_i(v,x)=\sum_{a=1}^3 c_a(v)w_i^{(a)}(x),\qquad
 \sum_i\partial_iC_i=0,\qquad
 C_j|_{B_r}=c_j.
 \tag{CE3}
\]
Every derivative of \(\chi\) in these formulas is retained.

The quaternionic bracket and matrix comparison remain
\[
 [u,v]=2u\times v,\quad
 T_a=-i\sigma_a/2,\quad [T_a,T_b]=\epsilon_{abc}T_c,\quad
 \varphi(e_a)=2T_a,\quad
 \langle U,V\rangle_{\mathrm{mat}}=-2\operatorname{tr}(UV).
 \tag{CE4}
\]
We use \(s=x^0=ct\), \(c>0\), the metric
\(\operatorname{diag}(-1,1,1,1)\), and coupling \(g>0\).
Time derivatives below are with respect to \(s\); in physical time they
are \(c^{-1}\partial_t\). In quaternionic coordinates the initial data are
\[
 A_0=0,\qquad A_i(0,x)=C_i(v,x),\qquad
 E_i(0,x)=\partial_s A_i(0,x)=0.
 \tag{CE5}
\]
The matrix connection is always \(\varphi(A)\); its energy has the full
coupling and trace factors in Section 4.

\subsection*{2. The established global input and its exact application}

Sung-Jin Oh's
\href{https://arxiv.org/abs/1210.1557v2}{global theorem, arXiv:1210.1557v2},
the subsection “Statement of the Main Theorem”, proves global existence and uniqueness in
temporal gauge for admissible data
\(\overline A_i\in\dot H^1(\mathbb R^3)\cap L^3(\mathbb R^3)\),
\(\overline E_i\in L^2(\mathbb R^3)\) satisfying
\(\sum_i(\partial_i\overline E_i+
[\overline A_i,\overline E_i])=0\).
Smooth data with all indicated Sobolev derivatives give a smooth solution;
its derivatives retain the stated Sobolev regularity for every finite time.
The companion
\href{https://arxiv.org/abs/1210.1558v2}{local theorem, arXiv:1210.1558v2},
Section 1.2, supplies uniqueness, difference estimates and persistence
of regularity. These are established analytical inputs, not hypotheses
postulated for this programme.

The original author TeX for both exact versions has been obtained and
preserved. In the global source the definitions and Main Theorem are at
lines 783–850, the continuation argument at 1030–1057, and the dependency
table at 2310–2373. The full heat-flow estimates are part of that cited
proof. We do not represent this application as a new proof of that
global PDE theorem or as an independent audit of its entire proof.

The convention map is explicit. Oh's time coordinate is our \(s=ct\),
and his Lie-algebra-valued connection is \(\varphi(A)\), with
\(SU(2)\) and the same metric signature and curvature sign. For the
matrix inner product displayed in his Section 1.1,
\[
 (U,V)_{\mathrm{Oh}}=\operatorname{tr}(UV^\star)
 =-\operatorname{tr}(UV)=\tfrac12\langle U,V\rangle_{\mathrm{mat}},
 \qquad
 (\varphi(u),\varphi(v))_{\mathrm{Oh}}=2\langle u,v\rangle_{\mathbb H}.
 \tag{CE6}
\]
Consequently, if \(\mathbf E_{\mathrm{Oh}}\) denotes his conserved
energy, our physical energy is \(2\mathbf E_{\mathrm{Oh}}/g^2\).
No coupling, generator or time factor is discarded.

Theorem 2.1. For every \(v\in W\), (CE5) has a unique global smooth
temporal-gauge solution of the source-free four-dimensional Yang–Mills
equations. Its spatial derivatives of every order and its time
derivatives have finite Sobolev norms on each finite time interval.
It has finite conserved physical energy. The assertion includes every
initial vertical rank \(0,3,6,9\) of the moment map.

Proof. Each component of \(C\) is a finite sum of a constant
quaternion \(c_a(v)\) times a fixed smooth compactly supported real
function. Thus \(C\in H^m\cap L^3\) for every nonnegative integer \(m\).
For example, with the usual Euclidean tuple norm,
\[
 \|C\|_{H^m}\leq\sum_a |c_a(v)|\,\|w^{(a)}\|_{H^m},\qquad
 \|C\|_{L^3}\leq\sum_a |c_a(v)|\,\|w^{(a)}\|_{L^3}.
 \tag{CE7}
\]
Every norm on the right is finite for the original \(\chi,r,R\).
Since the initial electric field is identically zero, its constraint
holds pointwise, including the cutoff region. Applying \(\varphi\)
preserves brackets and derivatives and multiplies squared norms by
the factors in (CE6). These are exactly the hypotheses of the cited
Main Theorem. Its existence, uniqueness and persistence statements
therefore apply to these original data without any smallness
restriction. Applying \(\varphi^{-1}\) returns the asserted solution.
The potential itself also lies in \(L^2\) on every finite time interval:
\[
 \|A(s)\|_{L^2}\leq\|C\|_{L^2}
       +|s|\sup_{|\sigma|\leq|s|}\|E(\sigma)\|_{L^2},
\]
by the temporal equation. Thus no missing zeroth-order Sobolev term is
inferred from a homogeneous norm.
Energy conservation is also proved directly in Section 3 below.
No coordinate in (CE1) was divided by, so zero and lower-rank inputs
are included. \(\square\)

This application replaces the previously unfinished classical
compact-support Cauchy step. It does not replace a quantum
construction by a classical solution.

\subsection*{3. Constraint propagation, finite propagation and the actual core}

Put
\[
 B_1=F_{23},\quad B_2=F_{31},\quad B_3=F_{12},\qquad
 (\operatorname{curl}_A U)_i
 =\sum_{j,k}\epsilon_{ijk}(\partial_jU_k+[A_j,U_k]).
\]
All sums in this section run from 1 to 3. With the conventions (CE4),
the temporal-gauge equations and Bianchi identity give
\[
 \partial_sA_i=E_i,\qquad
 \partial_s E=-\operatorname{curl}_A B,\qquad
 \partial_s B=\operatorname{curl}_A E,\qquad
 \sum_iD_iE_i=0.
 \tag{CE8}
\]
Indeed \(F_{ij}=\epsilon_{ijk}B_k\), so
\(\partial_sE_j=\sum_iD_iF_{ij}=-(\operatorname{curl}_AB)_j\).
Differentiating the definition of \(F_{ij}\) gives the equation for \(B\).
Conversely, along (CE8) the difference between \(B\) and
\[
 (\operatorname{curl}A)_i+\tfrac12
             \sum_{j,k}\epsilon_{ijk}[A_j,A_k]
\]
has zero time derivative. It vanishes initially, hence vanishes always.
For \(\mathcal G=\sum_iD_iE_i\), differentiation gives
\[
 \partial_s\mathcal G
 =-\tfrac12\sum_{i,j,k}\epsilon_{ijk}[F_{ij},B_k]
 =-\sum_k[B_k,B_k]=0.
 \tag{CE9}
\]
Thus the electric constraint is preserved by the evolution, rather
than reimposed at each time.

Let
\[
 e=\tfrac12\sum_i(|E_i|^2+|B_i|^2),\qquad
 q_j=-\sum_{i,k}\epsilon_{jik}\langle E_i,B_k\rangle_{\mathbb H}.
 \tag{CE10}
\]
Ad-invariance gives
\(\langle[X,Y],Z\rangle=-\langle Y,[X,Z]\rangle\).
On differentiating \(e\) using (CE8), the two connection terms cancel
under this identity and exchange of indices. The remaining ordinary
derivative terms give
\[
 \partial_se+\sum_j\partial_jq_j=0,\qquad
 |q\cdot n|\leq \tfrac12(|E|^2+|B|^2)=e
 \quad(|n|=1).
 \tag{CE11}
\]
For the inequality, write \(q\cdot n\) as the pairing of the spatial
cross-product map \(n\times E\) with \(B\); its operator norm is at most
one for each colour. Cauchy–Schwarz and \(2ab\leq a^2+b^2\) give the
displayed bound. Integrating the identity over space proves conservation.
One can justify the absence of a boundary term before proving support:
use cutoffs with gradient at most \(C/N\) supported on the annulus
\(N<|x|<2N\); the spacetime \(L^1\) norm of \(|q|\) on a finite time
interval is finite because the solution has continuous finite energy.
The error tends to zero.

Proposition 3.1. The actual solution, including its connection
coefficients in the specified temporal gauge, obeys
\[
 \operatorname{supp}(A(s),E(s),B(s))
 \subseteq\{x:\operatorname{dist}(x,\operatorname{supp}C)\leq |s|\}
 \subseteq\overline B_{R+|s|}(0).
 \tag{CE12}
\]
In physical time the speed in this formula is \(c\).

Proof. Fix a point \((s_0,x_0)\) with \(s_0>0\) whose backward
light cone misses \(\operatorname{supp}C\) on its initial disk.
Integrate (CE11) on the shrinking balls
\(B_{s_0-s}(x_0)\). The derivative of their energy is
\(-\int_{\partial B}(q\cdot n+e)\), which is nonpositive by (CE11).
The initial energy is zero, so \(E=B=0\) throughout that open cone,
and at its apex by continuity. For a fixed exterior \(x\), the vertical
segment from \((0,x)\) to \((s,x)\) also remains exterior. Integrating
\(\partial_sA=E=0\) and using \(C(x)=0\) proves \(A(s,x)=0\).
Reverse time to prove the negative-time assertion. \(\square\)

We also need domain of dependence for comparison with a nonzero field,
not only propagation of zero. Here is a direct uniqueness estimate.
For two smooth temporal-gauge solutions put
\(a=A-A'\), \(e_1=E-E'\), \(b_1=B-B'\) and
\[
 h=\tfrac12(|a|^2+|e_1|^2+|b_1|^2),\qquad
 q^{\mathrm{diff}}_j
 =-\sum_{i,k}\epsilon_{jik}\langle(e_1)_i,(b_1)_k\rangle.
\]
Subtracting (CE8), with \(E,B\) denoting the first solution, gives
\[
 \begin{aligned}
 \partial_sa&=e_1,\\
 \partial_s(e_1)_i
 &=-\bigl(\operatorname{curl}_{A'}b_1\bigr)_i
       -\sum_{j,k}\epsilon_{ijk}[a_j,B_k],\\
 \partial_s(b_1)_i
 &=\bigl(\operatorname{curl}_{A'}e_1\bigr)_i
       +\sum_{j,k}\epsilon_{ijk}[a_j,E_k].
 \end{aligned}\tag{CE13}
\]
The \(A'\)-terms cancel in the energy identity. Since
\(|[u,v]|\leq2|u||v|\), and each spatial alternating sum contains
six nonzero ordered triples, the remaining terms give the explicit bound
\[
 \partial_sh+\operatorname{div}q^{\mathrm{diff}}
 \leq [1+12(|E|+|B|)]h,\qquad
 |q^{\mathrm{diff}}\cdot n|\leq h.
 \tag{CE14}
\]
For example the term involving \(B\) has absolute value at most
\(12|e_1||a||B|\leq6|B|(|e_1|^2+|a|^2)\leq12|B|h\);
the \(E\)-term is identical, and
\(|\langle a,e_1\rangle|\leq h\).
Integrating on a shrinking cone and applying the integral form of
Gronwall's inequality proves equality throughout the cone when the
two initial triples agree on its initial disk. This argument is local
and therefore also compares the finite-energy solution to the smooth
homogeneous solution, whose total energy on \(\mathbb R^3\) is infinite.

Corollary 3.2. Let \(V_j\) be the exact homogeneous solution of
HC22 in the higher-carrier proof, with \(V_j(0)=c_j\).
Then the solution of (CE5) satisfies
\[
 A_j(s,x)=V_j(s)\quad\hbox{whenever}\quad |x|+|s|<r.
 \tag{CE15}
\]
The equality includes all spatial and time derivatives there.

Proof. The initial connection, electric field and curvature agree
on \(B_r\), by (CE3). Apply (CE13)–(CE14) to every backward or forward
cone with initial disk inside \(B_r\). Smooth equality on this open
diamond implies equality of derivatives. \(\square\)

For the retained equal-colour input \(c_j=e_j\), HC25 gives
\(V_j=f(s)e_j\), \(f''+8f^3=0\), \(f(0)=1\), \(f'(0)=0\), and
\(f'^2+4f^4=4\). Inside the diamond, in matrix coordinates,
\[
 F_{0j}=2f'T_j,\quad
 (F_{12},F_{23},F_{31})=4f^2(T_3,T_1,T_2).
 \tag{CE16}
\]
There is a useful gauge-invariant strengthening of the earlier magnetic
nonvanishing statement. Define the squared norms using (CE4), and keep
the indicated spacetime components fixed:
\[
 \begin{aligned}
 \mathcal I(s,x)
 &=|[F_{12},F_{23}]|_{\mathrm{mat}}^2
     +|[F_{01},F_{02}]|_{\mathrm{mat}}^2\\
 &=256f^8+16(f')^4
   =256\bigl((f^4)^2+(1-f^4)^2\bigr)\geq128.
 \end{aligned}\tag{CE17}
\]
The first commutator is \(16f^4T_2\), and the second is
\(4(f')^2T_3\). Since \(0\leq f^4\leq1\), the final inequality follows
from \(z^2+(1-z)^2=2(z-\tfrac12)^2+\tfrac12\).
Conjugation preserves both trace norms. Thus even at a zero of the
magnetic part the field in this diamond remains non-Abelian through
its electric curvature. This assertion is not a Lorentz-scalar claim;
it is a gauge-invariant statement about the stated components.

\subsection*{4. Full initial energy, and its exact zero set}

Use ordered pairs \(i<j\) and \(a<b\), retaining every summand:
\[
 H^a_{ij}=\partial_iw^{(a)}_j-\partial_jw^{(a)}_i,\qquad
 W^{ab}_{ij}=w^{(a)}_iw^{(b)}_j-w^{(b)}_iw^{(a)}_j.
 \tag{CE18}
\]
The nine derivative coefficients, with \(\chi_{ij}=\partial_i\partial_j\chi\),
are
\[
 \begin{array}{c|ccc}
 &H^1&H^2&H^3\\ \hline
 12&-2\chi_2-x_2(\chi_{11}+\chi_{22})
    &2\chi_1+x_1(\chi_{11}+\chi_{22})&x_1\chi_{23}\\
 13&-\chi_3-x_2\chi_{23}&x_1\chi_{23}
    &2\chi_1+x_1(\chi_{11}+\chi_{33})\\
 23&x_2\chi_{13}&-\chi_3-x_1\chi_{13}&\chi_2+x_1\chi_{12}.
 \end{array}\tag{CE19}
\]
All minors \(W^{ab}_{ij}\) are the exact two products specified by
(CE2) and (CE18); none is discarded. The curvature is
\[
 F_{ij}(C)=\sum_a c_aH^a_{ij}
             +\sum_{a<b}[c_a,c_b]W^{ab}_{ij}.
 \tag{CE20}
\]
Define the finite, cutoff-dependent real coefficients
\[
 \begin{aligned}
 I_{ab}&=\int_{\mathbb R^3}\sum_{i<j}H^a_{ij}H^b_{ij}\,d^3x,\\
 J_{a,bc}&=\int_{\mathbb R^3}\sum_{i<j}H^a_{ij}W^{bc}_{ij}\,d^3x
                   &&(b<c),\\
 K_{ab,cd}&=\int_{\mathbb R^3}\sum_{i<j}W^{ab}_{ij}W^{cd}_{ij}\,d^3x
                   &&(a<b,\ c<d).
 \end{aligned}\tag{CE21}
\]
Their complete contribution to the conserved physical energy is
\[
 \begin{aligned}
 \mathscr E(v)
 &=\frac1{2g^2}\int\left(
       \sum_i|\varphi(E_i)|_{\mathrm{mat}}^2+
       \sum_{i<j}|\varphi(F_{ij})|_{\mathrm{mat}}^2\right)d^3x\\
 &=\frac{4}{g^2}\int e\,d^3x\\
 &=\frac2{g^2}\left[
     \sum_{a,b}\langle c_a,c_b\rangle I_{ab}
     +2\sum_{a,\ b<c}\langle c_a,[c_b,c_c]\rangle J_{a,bc}
     +\sum_{a<b,\ c<d}\langle[c_a,c_b],[c_c,c_d]\rangle K_{ab,cd}
              \right].
 \end{aligned}\tag{CE22}
\]
The last equality evaluates the first line at \(s=0\), where \(E=0\),
and expands (CE20). In the original line coordinates it contains
degrees four, six and eight. No radial or parity condition on \(\chi\)
has been assumed, so in particular the mixed \(J\)-terms are retained.
For the equal-colour input the original ball contributes
\[
 \mathscr E(v)\geq\frac{24}{g^2}\operatorname{vol}(B_r)
               =\frac{32\pi r^3}{g^2}>0.
 \tag{CE23}
\]
This is a lower bound; the annulus energy in (CE22) is still present.

Proposition 4.1. The zero set of \(\mathscr E\) in the complete
rank-24 fibre is exactly the retained rank-12 linear subspace
\(\alpha=\beta=\gamma=0\). This also describes exactly the inputs
whose entire evolved field is zero.

Proof. If the three indicated line coordinates vanish, \(C=0\)
and uniqueness gives the zero solution. Conversely, \(\mathscr E=0\)
implies \(F_{ij}(C)=0\) by positivity and smoothness.
We prove that a smooth compactly supported divergence-free flat
quaternionic connection \(C\) is zero, with no gauge choice added.
Set
\[
 T_{ij}=\langle C_i,C_j\rangle-\tfrac12\delta_{ij}\sum_k|C_k|^2.
\]
Using \(\sum_i\partial_iC_i=0\) and flatness,
\[
 \begin{aligned}
 \sum_i\partial_iT_{ij}
 &=\sum_i\langle C_i,\partial_iC_j-\partial_jC_i\rangle\\
 &=-\sum_i\langle C_i,[C_i,C_j]\rangle=0.
 \end{aligned}
\]
Integrate the divergence of \(\sum_jx_jT_{ij}\) over \(\mathbb R^3\).
Its compact support removes the boundary term, while its trace is
\(\sum_iT_{ii}=-\tfrac12\sum_i|C_i|^2\). Hence \(C=0\).
On \(B_r\), (CE3) then gives all \(c_j=0\). The exact identity
\(|a e\bar a|=|a|^2\) for \(|e|=1\) gives
\(\alpha=\beta=\gamma=0\). The remaining twelve real coordinates are
unrestricted. Finally a solution vanishing identically has zero initial
data, and the converse was already proved. \(\square\)

This also covers nonzero inputs whose colours commute: their cutoff
energy is positive even when the constant core's magnetic energy is zero.
When all \(c_a\) lie in one fixed quaternionic line, that line is
bracket-abelian. The linear Maxwell solution within it solves (CE8);
uniqueness identifies it with the full solution. Such inputs are
explicitly included, not incorrectly described as interacting.

\subsection*{5. Smooth parameter dependence, with a constructive derivative equation}

We prove the additional parameter statement needed for bundle descent.
It is stronger than merely assigning a solution separately to each input.
Work with the complete real tuple \(U=(A,E,B)\). Define
\[
 \mathcal L(A,E,B)=(0,-\operatorname{curl}B,\operatorname{curl}E),
 \quad N(U)=N_1(U)+Q(U,U),
 \tag{CE24}
\]
where \(N_1(U)=(E,0,0)\), and the bilinear map, with ordered arguments,
is
\[
 Q(U,V)=\left(
 0,\ -\sum_{j,k}\epsilon_{ijk}[A_j(U),B_k(V)],\
       \sum_{j,k}\epsilon_{ijk}[A_j(U),E_k(V)]
 \right)_{i=1}^3.
 \tag{CE25}
\]
There are no spatial derivatives in \(N\). Equation (CE8) is exactly
\[
 U(s)=e^{s\mathcal L}U(0)
           +\int_0^s e^{(s-\sigma)\mathcal L}N(U(\sigma))\,d\sigma.
 \tag{CE26}
\]

Here are explicit estimates sufficient to justify all parameter
differentiations. For an integer \(m\geq2\), use
\[
 \widehat f(\xi)=\int e^{-ix\cdot\xi}f(x)\,d^3x,\qquad
 \|f\|_{H^m}^2=(2\pi)^{-3}\int(1+|\xi|^2)^m|\widehat f(\xi)|^2\,d^3\xi.
\]
For tuples take the sum of the squared component norms.
The Fourier symbol of \(\mathcal L\) is skew-Hermitian, since the curl
operator is self-adjoint in \(L^2\). Thus \(e^{s\mathcal L}\) is an
isometry on every such \(H^m\).
The convolution formula for a product, the inequality
\(\langle\xi\rangle^m\leq2^{m-1}
(\langle\eta\rangle^m+\langle\xi-\eta\rangle^m)\), and
Cauchy–Schwarz give
\[
 \|fg\|_{H^m}\leq C_m\|f\|_{H^m}\|g\|_{H^m},\qquad
 C_m=2^m(2\pi)^{-3/2}
 \left(\frac{\pi^{3/2}\Gamma(m-\tfrac32)}{\Gamma(m)}\right)^{1/2}.
 \tag{CE27}
\]
To check the constant, the \(L^1\) norm of \(\widehat f\) is at most
\((\int\langle\xi\rangle^{-2m}d^3\xi)^{1/2}
 \|\langle\xi\rangle^m\widehat f\|_2\).
The integral equals the Gamma quotient in (CE27), by polar coordinates
and the substitution \(u=|\xi|^2\). Applying Young's inequality to
each of the two convolution terms gives precisely the stated factor.

In each of the last two blocks of (CE25), expansion of the spatial
and colour alternating tensors gives 36 nonzero ordered products,
each with coefficient of absolute value 2. Bounding the norm of a
sum by the sum of its norms therefore proves
\[
 \|Q(U,V)\|_{H^m}\leq144C_m\|U\|_{H^m}\|V\|_{H^m}.
 \tag{CE28}
\]
The estimate retains both blocks; it is an upper bound, not an equality.
Consequently on the ball \(\|U\|,\|V\|\leq2M\),
\[
 \|N(U)-N(V)\|_{H^m}
 \leq(1+576C_mM)\|U-V\|_{H^m}.
 \tag{CE29}
\]
For \(M\geq\max(1,\|U(0)\|_{H^m})\) and
\[
 d_m(M)=\frac1{4(1+576C_mM)},
 \tag{CE30}
\]
the right side of (CE26) maps the closed radius-\(2M\) ball of
\(C([-d_m,d_m],H^m)\) into itself and has contraction constant at most
\(1/4\). Indeed its size is at most
\(M+d_m(2M+576C_mM^2)\leq2M\), and (CE29) proves contraction.
This constructs the local solution and its uniqueness in that space.
The curvature and constraint identities (CE9) show that it is the
same Yang–Mills solution as Theorem 2.1 when the data have the required
constraints.

Let \(y\) denote any finite collection of real parameter coordinates.
The initial triple
\[
 U(0;y)=(C(y),\,0,\,(F_{23}(C(y)),F_{31}(C(y)),F_{12}(C(y))))
 \tag{CE31}
\]
is a polynomial map into every \(H^m\) in a local bundle frame.
For a nonzero parameter multiindex \(\beta\), its derivative satisfies
the following complete equation:
\[
 \begin{aligned}
 U_\beta(s)
 &=e^{s\mathcal L}\partial_y^\beta U(0)
 +\int_0^s e^{(s-\sigma)\mathcal L}\biggl(
 N_1(U_\beta)+Q(U_\beta,U)+Q(U,U_\beta)\\
 &\hspace{32mm}
 +\sum_{\substack{0<\gamma<\beta\\\gamma\leq\beta}}
        {\beta\choose\gamma}Q(U_\gamma,U_{\beta-\gamma})
                    \biggr)(\sigma)\,d\sigma.
 \end{aligned}\tag{CE32}
\]
In the sum, \(0<\gamma<\beta\) means that neither \(\gamma=0\) nor
\(\gamma=\beta\); comparisons are componentwise.
No differentiated term is removed.

For completeness, the differentiability assertion follows from the
same contraction rather than a formal differentiation alone.
Subtract the fixed-point equations at \(y+h\) and \(y\).
Estimate (CE29) first gives continuity and a uniform difference bound.
After division by a coordinate increment, the error from replacing
the difference quotient by the solution of the first-derivative version
of (CE32) is bounded by a constant times the product of the data
increment and the solution increment. This tends to zero by (CE28);
the inverse of the linear fixed-point operator has norm at most
\((1-\tfrac14)^{-1}\). This proves the first derivative.
Repeated difference quotients give (CE32): the product rule for the
bilinear \(Q\) gives exactly the displayed binomial sum, whose lower
derivatives are already continuous by induction. The same inverse
bound proves existence and continuity of the next derivative.
This proves smooth dependence on every finite parameter collection.

For a fixed input \(y_0\) and finite \(T\), the global smooth solution
from Theorem 2.1 has finite
\(M_T=1+\sup_{|s|\leq T}\|U(s;y_0)\|_{H^m}\).
Applying the preceding construction successively on finitely many
intervals of length at most \(d_m(2M_T)\) proves smooth dependence
near \(y_0\) throughout \([-T,T]\). Local uniqueness identifies each
extension with the global solution already constructed.
Covering a compact parameter set by finitely many such neighbourhoods
shows that every parameter derivative is bounded in \(H^m\) there,
uniformly on \([-T,T]\). This argument makes no bound uniform as a
parameter tends to infinity. Applying it for every \(m\), followed by
Sobolev embedding and (CE8), proves joint smoothness in parameters,
space and time.

\subsection*{6. Descent over both original bases and all mixed curvatures}

Simultaneous conjugation by \(u\in\operatorname{Sp}(1)\) sends
\(C(v)\) to \(uC(v)\bar u=C(uv)\). It preserves (CE8) and its zero
initial electric field. Uniqueness therefore gives
\[
 A(uv;s,x)=uA(v;s,x)\bar u.
 \tag{CE33}
\]
Let \(\mathscr S\) be the set of smooth temporal-gauge solutions
constructed above, equipped with the topology of uniform convergence
on compact time intervals in every spatial \(H^m\), and all time
derivatives. The smooth solution map induces
\[
 \mathfrak E_D:E_D=L\times_K W\longrightarrow L\times_K\mathscr S,
 \qquad [\ell,v]\longmapsto[\ell,A(v)],
 \quad
 r_H^*\mathfrak E_D=\mathfrak E_H.
 \tag{CE34}
\]
This is well defined by (CE33). Its initial-time restriction is exactly
\(\mathcal M_D\), respectively \(\mathcal M_H\). This proves the
connecting morphism while preserving all initial rank strata and the
rank-12 zero subbundle. No global frame for the nontrivial colour
bundle, or nowhere-zero global section, has been selected.

There is also an exact differential statement. Let \(\operatorname{ev}_0\)
evaluate the connection at \(s=0\), and let \(\mathcal T\) denote the
Cauchy solution map on the finite-dimensional space of coefficients
\((c_1,c_2,c_3)\) in (CE3). Then
\(\mathfrak E=\mathcal T\mathcal M\) and
\(\operatorname{ev}_0\mathfrak E=\mathcal M\). The chain rule and these
two factorizations imply
\[
 \operatorname{rank}d_v\mathfrak E
   =\operatorname{rank}d_v\mathcal M
   =3\,\#\{a\in\{\alpha,\beta,\gamma\}:a\ne0\}.
 \tag{CE34a}
\]
Indeed the first factorization gives the upper bound, and applying
\(d\operatorname{ev}_0\) gives the reverse bound. In particular the
solution map creates no new independent input directions.
Two inputs give the same entire solution precisely when they give the
same three moment-map values. The free spectator coordinates remain
irrelevant to the solution. For a nonzero line coordinate \(a\) with
label \(e\), all coordinates with the same value \(ae\bar a\) are
exactly \(a(\cos\theta+e\sin\theta)\). To prove this, equality of moment
maps first gives equal norms; writing \(b=au\) therefore gives
\(|u|=1\) and \(ue\bar u=e\). The imaginary part of a quaternion
commuting with \(e\) is parallel to \(e\), so precisely the displayed
circle is obtained. The fibre over zero consists only of \(a=0\).
This exact redundancy must be retained when later computing Gram
kernels; distinct coordinate labels alone do not create distinct states.

Set \(Y=\operatorname{Tot}(E_D)\), as before. The tautological input
and the retained Maurer–Cartan base connection \(b=b_A\,dy^A\)
give the full connection
\[
 \mathbb B=b_A\,dy^A+\varphi(A_i(y;s,x))\,dx^i,\qquad \mathbb B_0=0.
 \tag{CE35}
\]
Its local components transform by the full connection law, because
frame changes depend on \(y\) and (CE33) controls the physical components.
Its curvature is
\[
 \begin{aligned}
 \mathbb F_{AB}&=\partial_Ab_B-\partial_Bb_A+[b_A,b_B],\\
 \mathbb F_{Ai}&=\partial_A\varphi(A_i)+[b_A,\varphi(A_i)],\\
 \mathbb F_{ij}&=\varphi(\partial_iA_j-\partial_jA_i+[A_i,A_j]),\\
 \mathbb F_{A0}&=0,\qquad \mathbb F_{0i}=\varphi(E_i).
 \end{aligned}\tag{CE36}
\]
The mixed parameter derivatives exist with the compact-parameter bounds
proved in Section 5. In a quaternionic local frame write
\(h_A=\varphi^{-1}(b_A)\) and
\(Z_{A i}=\partial_AA_i+[h_A,A_i]\).
Then \(Z_{A i}\) obeys the exact linearized physical equations
\[
 \begin{aligned}
 \partial_s^2 Z_{A j}
 &=\sum_i\bigl([Z_{A i},F_{ij}]
       +D_i(D_iZ_{A j}-D_jZ_{A i})\bigr),\\
 0&=\sum_i\bigl(D_i\partial_sZ_{A i}+[Z_{A i},E_i]\bigr),\\
 Z_{A i}(0)&=\partial_AC_i+[h_A,C_i],\qquad
 \partial_sZ_{A i}(0)=0.
 \end{aligned}\tag{CE37}
\]
To prove this, differentiate the physical Yang–Mills equation and Gauss
constraint covariantly in the parameter. The base form is independent
of \(s,x\), so differentiation of \(D_iF_{ij}\) gives exactly the two
displayed terms; differentiation of \(F_{ij}\) gives
\(D_iZ_{A j}-D_jZ_{A i}\). Formula (CE32) proves these derivatives
exist, and so justifies this calculation. The base curvature remains
the one computed in HC11 and need not vanish outside the spatial
support. No equation in the 73 parameter dimensions is being added
to the four-dimensional physical Yang–Mills system.

\subsection*{7. Exact dilation, including its moving support and energy}

This construction now permits an actual all-time finite-energy family,
with every scale retained. For \(\lambda>0\), define
\[
 \chi_\lambda(x)=\chi(x/\lambda),\qquad
 w^{(a)}_\lambda(x)=w^{(a)}(x/\lambda),\qquad
 v_\lambda=\lambda^{-1/2}v.
\]
The second identity follows directly by taking the original three
curls with \(\chi_\lambda x_2,-\chi_\lambda x_1,\chi_\lambda x_1\);
the derivative of the coordinate factor cancels the corresponding
chain-rule factor, leaving exactly (CE2) at \(x/\lambda\).
The inner and outer radii are now \(\lambda r,\lambda R\).
Since \(\mu_e(\lambda^{-1/2}a)=\lambda^{-1}\mu_e(a)\),
\[
 C^{[\lambda]}(v,x)
 =\sum_a c_a(v_\lambda)w^{(a)}_\lambda(x)
 =\lambda^{-1}C(v,x/\lambda).
 \tag{CE38}
\]
The brackets, time derivatives and spatial derivatives then give
\[
 \begin{aligned}
 A^{[\lambda]}(s,x)&=\lambda^{-1}A(s/\lambda,x/\lambda),\\
 E^{[\lambda]}(s,x)&=\lambda^{-2}E(s/\lambda,x/\lambda),\\
 F_{ij}^{[\lambda]}(s,x)&=\lambda^{-2}F_{ij}(s/\lambda,x/\lambda),\\
 D_{A^{[\lambda]}}^\mu F_{\mu\nu}^{[\lambda]}
   &=\lambda^{-3}(D_A^\mu F_{\mu\nu})(s/\lambda,x/\lambda)=0,\\
 \mathscr E^{[\lambda]}&=\lambda^{-1}\mathscr E,\qquad
 \operatorname{supp}A^{[\lambda]}(s)
       \subseteq\overline B_{\lambda R+|s|}.
 \end{aligned}\tag{CE39}
\]
The first field has exactly the Cauchy data (CE38) and zero electric
data, so uniqueness identifies it with the solution already constructed
from the scaled cutoff and line input. There is no new approximation.
The energy factor is the product \(\lambda^{-4}\lambda^3\) from curvature
and volume. Also
\[
 \|C^{[\lambda]}\|_{L^2}=\lambda^{1/2}\|C\|_{L^2},\quad
 \|\nabla C^{[\lambda]}\|_{L^2}=\lambda^{-1/2}\|\nabla C\|_{L^2},\quad
 \|C^{[\lambda]}\|_{L^3}=\|C\|_{L^3}.
 \tag{CE40}
\]
These recover the source theorem's scaling convention without
replacing the programme field by a unit-scale field.
For the equal-colour input the exact core diamond is now
\(|x|+|s|<\lambda r\), and (CE17) becomes
\[
 \mathcal I^{[\lambda]}(s,x)
 =\lambda^{-8}\mathcal I(s/\lambda,x/\lambda)
 \geq128\lambda^{-8}.
 \tag{CE41}
\]
Thus each finite scale has a smooth finite-energy non-Abelian region.
Its energy tends to zero and its support scale grows as
\(\lambda\to\infty\). Those are classical facts about this explicit
family, not a positive-energy spectral-weight assertion for a quantum
Hamiltonian. In particular the simultaneous decrease of amplitude is
retained; this calculation cannot be used to assume that a quantum
excitation remains orthogonal to the vacuum.

\subsection*{8. The exact receiving map to link holonomies}

To specify the next quantum calculation on the fields just constructed,
fix \(s\), any finite oriented graph embedded by piecewise smooth
spatial paths \(\gamma_e:[0,1]\to\mathbb R^3\), and use the matrix
connection \(\mathcal A=\varphi(A)\). Define
\[
 \frac{d}{d\tau}U_e(\tau)
 =-\sum_i\mathcal A_i(s,\gamma_e(\tau))\dot\gamma_e^i(\tau)U_e(\tau),
 \quad U_e(0)=I,\qquad H_e=U_e(1).
 \tag{CE42}
\]
The coefficients are smooth on each path piece. Picard iteration of
the integral equation gives its unique solution; successive pieces
are joined by their endpoint values. Differentiation gives
\(\partial_\tau(U_e^\star U_e)=0\), and the trace of the coefficient
is zero, so \(\det U_e=1\). Therefore \(H_e\in SU(2)\).
Reversing the path gives \(H_{\bar e}=H_e^{-1}\).

For the gauge convention
\(\mathcal A^q=q\mathcal A q^{-1}-dq\,q^{-1}\), direct substitution
in (CE42) gives
\[
 H_e(\mathcal A^q)
   =q(\gamma_e(1))H_e(\mathcal A)q(\gamma_e(0))^{-1}.
 \tag{CE43}
\]
For a variation \(z\) of the field, differentiation of its integral
equation gives the complete insertion formula
\[
 \delta H_e
 =-H_e\int_0^1 U_e(\tau)^{-1}
       \left(\sum_i\varphi(z_i)(s,\gamma_e(\tau))
                          \dot\gamma_e^i(\tau)\right)
       U_e(\tau)\,d\tau.
 \tag{CE44}
\]
It follows either by differentiating \(U_e^{-1}\delta U_e\), or by
variation of constants; both retain the order of the matrix factors.
The bounds of Section 5 on any compact path set justify differentiation
for all fibre parameters. Consequently (CE42) is a smooth,
gauge-equivariant map from the actual evolved bundle family to the
edge-configuration space, including every input rank stratum.

There is also an exact relation to dilation. For
\(\gamma_e^{[\lambda]}=\lambda\gamma_e\), the product of the
\(\lambda^{-1}\) connection factor and the \(\lambda\) tangent factor
in (CE42) gives
\[
 H_{\lambda\gamma_e}\bigl(A^{[\lambda]}(s)\bigr)
        =H_{\gamma_e}\bigl(A(s/\lambda)\bigr).
 \tag{CE45}
\]
Holding an embedded graph fixed is a different limit and does not
follow by erasing the path factor in this formula.

The next construction is the physical gauge projection of vectors
centered at these actual edge configurations, with their full Gram
and Hamiltonian kernels. It must also receive the retained period
matrix and angular monodromy. The spectral weight, vacuum subtraction,
fixed regulator path and continuum reconstruction remain to be proved.
The classical dilation formulas alone supply none of those assertions.

\subsection*{9. Sources and reproducibility}

The original moment-map and higher-carrier arguments are linked in
Section 1, with their public proof locators. The external global
existence theorem and its local companion are Sung-Jin Oh's exact
arXiv v2 author sources, dated 29 June 2015. Their archives, hashes,
reading coverage and exact receiving equations are recorded in
SOURCE\_PROVENANCE.md and the private source-reading ledger.

The algebraic checker checks the cutoff curls and derivatives,
constraint identities, full initial energy expansion, stress identity,
core commutators and scale factors. These finite checks do not certify
the imported analytical existence theorem. The full local estimates,
source application, connecting maps and all new arguments are written
above.

\textit{The reproducible figure and inspected rendering are retained with the corresponding Markdown proof.}

Figure: Proposition 3.1 and Corollary 3.2 retain the expanding support
bound and the exact inner diamond. The right panel renders the ODE
of (CE16), with its two curvature commutator norms and the proved
bound (CE17). The spatial example uses \(r=1,R=2\) only for drawing;
the proof retains arbitrary \(0<r<R\). The plotted ODE is a numerical
rendering of an exact analytical construction. The source is
figures/compact\_cauchy\_figure.py.

\subsection*{Finite physical-state continuation, 8 October 2026}

The \href{https://github.com/KokunoYumeto/yang-mills-interacting-workbench/blob/main/yang-mills/continuations/20260930-s6-ns-moment-map-bridge/PERIOD_COUPLING_AND_PHYSICAL_KERNELS.md}{complete period and physical-kernel proof} now supplies the
full angular pullback, exact physical gauge projection, actual Gram and
Hamiltonian kernels, and a nonzero vector orthogonal to the interacting
finite-lattice vacuum. Equations PK43–PK43g correct the joint path while
retaining the same coupling and every original cutoff energy term.

Theorem 9.1 identifies the complete first electric layer. Equations PK47–PK49
strengthen the raw norm using all faces in both reflected core cubes.
Equations PK50–PK54 retain the full magnetic operator, calculate its first
compression and the exact complementary resolvent return. The continuum
spectral measure, reconstruction and theory identification remain the active
calculation; decreasing classical energy does not settle them.
\end{document}
